% --- Archon physics formalization source begin ---
% archon:physics
% archon:covers PhyXMiniProblems/problem_phyx_mini_0071.lean
% archon:source-report reports/phyx_mini/problem_phyx_mini_0071.source.json

\chapter{Physics problem phyx\_mini\_0071}
\label{ch:PhyXMiniProblems_problem_phyx_mini_0071}

\paragraph{Problem source.}
\#\# Physical scenario

The 80-cm-tall, 65-cm-wide tank shown in Figure is completely filled with water. The tank has marks every 10 cm along one wall, and the 0 cm mark is barely submerged. As you stand beside the opposite wall, your eye is level with the top of the water.

\#\# Question

What mark do you see on the meter stick if the tank is empty?

\#\# Answer choices

- A: A: 20cm

- B: B: 50cm

- C: C: 60cm

- D: D: 80cm

\#\# Figure caption (auxiliary; use the image as primary evidence)

The image is a diagram of a rectangular container filled with a light blue substance, possibly representing water. The height of the container is marked on the right-hand side with a depth scale labeled "Depth (cm)" that ranges from 0 cm at the top to 80 cm at the bottom in increments of 10 cm.

At the top left corner, there is an icon of an eye labeled "Observation point," indicating the viewpoint from which the depth is being observed. 

The width of the container at the bottom is marked as "65 cm" with a double-headed arrow indicating the measurement.

Overall, the diagram appears to illustrate the dimensions and observation point for examining the depth within the container.

\#\# Dataset metadata

Geometrical Optics; Spatial Relation Reasoning, Physical Model Grounding Reasoning

\paragraph{Recorded answer/context.}
C: C: 60cm

\paragraph{Figure/image path.}
/root/proposal\_for\_physic/science-mango/phyx\_mini\_run/phyx\_data/test\_image/71.png

\paragraph{Formalization target.}
create a compiling Lean file with sorry bodies at `PhyXMiniProblems/problem\_phyx\_mini\_0071.lean`. The Lean declarations must preserve the physical quantities, dimensions or dimensional roles, figure labels, governing-law hypotheses, and final relation expressed by this problem.
Use Mathlib/Physlib names found through LeanExplore where available. If a physics API is missing, introduce faithful local abstractions rather than scalar placeholder aliases.

\begin{theorem}[Physics formalization target]
\label{thm:physics:phyx_mini_0071:target}
\lean{PhyXMiniProblems.ProblemPhyXMini0071.problem_phyx_mini_0071}
\uses{def:physics:phyx-mini-0071:phyxminiproblems-problemphyxmini0071-lengthincentimeters, def:physics:phyx-mini-0071:phyxminiproblems-problemphyxmini0071-tankviewingsetup, def:physics:phyx-mini-0071:phyxminiproblems-problemphyxmini0071-matchestankproblemandfigure, def:physics:phyx-mini-0071:phyxminiproblems-problemphyxmini0071-modelsordinarywaterandair, def:physics:phyx-mini-0071:phyxminiproblems-problemphyxmini0071-satisfiesgrazingwaterairrefraction, def:physics:phyx-mini-0071:phyxminiproblems-problemphyxmini0071-satisfieslimitingsightlinegeometry, def:physics:phyx-mini-0071:phyxminiproblems-problemphyxmini0071-iscorrectanswer}
The assigned autoformalize agent should translate this physics problem into Lean declarations in the covered file, with theorem and lemma proof bodies written as `by sorry`.
\end{theorem}
\begin{proof}
This is an autoformalization task, not a proof task. Produce statements that can later be proved by the physics prover without weakening the physical model.
\end{proof}

% --- Archon declaration topology begin ---
\section{Declaration topology}
\begin{definition}[Dim Length]
  \label{def:physics:phyx-mini-0071:phyxminiproblems-problemphyxmini0071-dimlength}
  \lean{PhyXMiniProblems.ProblemPhyXMini0071.DimLength}
  A physical length, represented independently of the unit used to read it
\end{definition}
\begin{proof}
  This is the declared model definition of Dim Length.
\end{proof}
\begin{definition}[centimeter Unit Choices]
  \label{def:physics:phyx-mini-0071:phyxminiproblems-problemphyxmini0071-centimeterunitchoices}
  \lean{PhyXMiniProblems.ProblemPhyXMini0071.centimeterUnitChoices}
  Unit choices in which length readouts are expressed in centimeters
\end{definition}
\begin{proof}
  This is the declared model definition of centimeter Unit Choices.
\end{proof}
\begin{definition}[length In Centimeters]
  \label{def:physics:phyx-mini-0071:phyxminiproblems-problemphyxmini0071-lengthincentimeters}
  \lean{PhyXMiniProblems.ProblemPhyXMini0071.lengthInCentimeters}
  \uses{def:physics:phyx-mini-0071:phyxminiproblems-problemphyxmini0071-dimlength, def:physics:phyx-mini-0071:phyxminiproblems-problemphyxmini0071-centimeterunitchoices}
  The scalar centimeter readout of a physical length
\end{definition}
\begin{proof}
  This is the declared model definition of length In Centimeters.
\end{proof}
\begin{definition}[Tank Wall]
  \label{def:physics:phyx-mini-0071:phyxminiproblems-problemphyxmini0071-tankwall}
  \lean{PhyXMiniProblems.ProblemPhyXMini0071.TankWall}
  The two opposite vertical walls in the cross-sectional ray diagram
\end{definition}
\begin{proof}
  This is the declared model definition of Tank Wall.
\end{proof}
\begin{definition}[Tank Contents]
  \label{def:physics:phyx-mini-0071:phyxminiproblems-problemphyxmini0071-tankcontents}
  \lean{PhyXMiniProblems.ProblemPhyXMini0071.TankContents}
  The physical filling states mentioned by the source wording and figure
\end{definition}
\begin{proof}
  This is the declared model definition of Tank Contents.
\end{proof}
\begin{definition}[source Question Contents]
  \label{def:physics:phyx-mini-0071:phyxminiproblems-problemphyxmini0071-sourcequestioncontents}
  \lean{PhyXMiniProblems.ProblemPhyXMini0071.sourceQuestionContents}
  \uses{def:physics:phyx-mini-0071:phyxminiproblems-problemphyxmini0071-tankcontents}
  The filling state named literally by the final sentence of the question
\end{definition}
\begin{proof}
  This is the declared model definition of source Question Contents.
\end{proof}
\begin{definition}[pictured Scenario Contents]
  \label{def:physics:phyx-mini-0071:phyxminiproblems-problemphyxmini0071-picturedscenariocontents}
  \lean{PhyXMiniProblems.ProblemPhyXMini0071.picturedScenarioContents}
  \uses{def:physics:phyx-mini-0071:phyxminiproblems-problemphyxmini0071-tankcontents}
  The filling state shown in the diagram and stated in the scenario
\end{definition}
\begin{proof}
  This is the declared model definition of pictured Scenario Contents.
\end{proof}
\begin{definition}[Optical Medium]
  \label{def:physics:phyx-mini-0071:phyxminiproblems-problemphyxmini0071-opticalmedium}
  \lean{PhyXMiniProblems.ProblemPhyXMini0071.OpticalMedium}
  The homogeneous optical media separated by the horizontal water surface
\end{definition}
\begin{proof}
  This is the declared model definition of Optical Medium.
\end{proof}
\begin{definition}[Depth Mark]
  \label{def:physics:phyx-mini-0071:phyxminiproblems-problemphyxmini0071-depthmark}
  \lean{PhyXMiniProblems.ProblemPhyXMini0071.DepthMark}
  Labels of all 10 cm depth marks drawn from the top to the tank bottom
\end{definition}
\begin{proof}
  This is the declared model definition of Depth Mark.
\end{proof}
\begin{definition}[depth Mark Centimeters]
  \label{def:physics:phyx-mini-0071:phyxminiproblems-problemphyxmini0071-depthmarkcentimeters}
  \lean{PhyXMiniProblems.ProblemPhyXMini0071.depthMarkCentimeters}
  \uses{def:physics:phyx-mini-0071:phyxminiproblems-problemphyxmini0071-depthmark}
  The centimeter depth printed beside each mark in the figure
\end{definition}
\begin{proof}
  This is the declared model definition of depth Mark Centimeters.
\end{proof}
\begin{definition}[Tank Depth Scale]
  \label{def:physics:phyx-mini-0071:phyxminiproblems-problemphyxmini0071-tankdepthscale}
  \lean{PhyXMiniProblems.ProblemPhyXMini0071.TankDepthScale}
  \uses{def:physics:phyx-mini-0071:phyxminiproblems-problemphyxmini0071-dimlength, def:physics:phyx-mini-0071:phyxminiproblems-problemphyxmini0071-tankwall}
  The depth scale fixed to one wall of the tank
\end{definition}
\begin{proof}
  This is the declared model definition of Tank Depth Scale.
\end{proof}
\begin{definition}[Observation Point]
  \label{def:physics:phyx-mini-0071:phyxminiproblems-problemphyxmini0071-observationpoint}
  \lean{PhyXMiniProblems.ProblemPhyXMini0071.ObservationPoint}
  \uses{def:physics:phyx-mini-0071:phyxminiproblems-problemphyxmini0071-dimlength, def:physics:phyx-mini-0071:phyxminiproblems-problemphyxmini0071-tankwall}
  The observer location shown at the upper rim of the opposite wall
\end{definition}
\begin{proof}
  This is the declared model definition of Observation Point.
\end{proof}
\begin{definition}[Tank Viewing Setup]
  \label{def:physics:phyx-mini-0071:phyxminiproblems-problemphyxmini0071-tankviewingsetup}
  \lean{PhyXMiniProblems.ProblemPhyXMini0071.TankViewingSetup}
  \uses{def:physics:phyx-mini-0071:phyxminiproblems-problemphyxmini0071-dimlength, def:physics:phyx-mini-0071:phyxminiproblems-problemphyxmini0071-tankcontents, def:physics:phyx-mini-0071:phyxminiproblems-problemphyxmini0071-opticalmedium, def:physics:phyx-mini-0071:phyxminiproblems-problemphyxmini0071-tankdepthscale, def:physics:phyx-mini-0071:phyxminiproblems-problemphyxmini0071-observationpoint}
  The unknown sightline depth is kept as a physical length. In particular, this structure contains no field fixing it to the requested 60 cm mark
\end{definition}
\begin{proof}
  This is the declared model definition of Tank Viewing Setup.
\end{proof}
\begin{definition}[Matches Tank Problem And Figure]
  \label{def:physics:phyx-mini-0071:phyxminiproblems-problemphyxmini0071-matchestankproblemandfigure}
  \lean{PhyXMiniProblems.ProblemPhyXMini0071.MatchesTankProblemAndFigure}
  \uses{def:physics:phyx-mini-0071:phyxminiproblems-problemphyxmini0071-lengthincentimeters, def:physics:phyx-mini-0071:phyxminiproblems-problemphyxmini0071-picturedscenariocontents, def:physics:phyx-mini-0071:phyxminiproblems-problemphyxmini0071-depthmarkcentimeters, def:physics:phyx-mini-0071:phyxminiproblems-problemphyxmini0071-tankviewingsetup}
  Numerical and qualitative data read from the problem and figure: an 80 cm deep, 65 cm wide tank filled to the top, a depth scale every 10 cm, and an eye at water level on the opposite wall. No limiting-ray depth is assigned
\end{definition}
\begin{proof}
  This is the declared model definition of Matches Tank Problem And Figure.
\end{proof}
\begin{definition}[Models Ordinary Water And Air]
  \label{def:physics:phyx-mini-0071:phyxminiproblems-problemphyxmini0071-modelsordinarywaterandair}
  \lean{PhyXMiniProblems.ProblemPhyXMini0071.ModelsOrdinaryWaterAndAir}
  \uses{def:physics:phyx-mini-0071:phyxminiproblems-problemphyxmini0071-tankviewingsetup}
  Standard refractive-index readouts used for ordinary water and air
\end{definition}
\begin{proof}
  This is the declared model definition of Models Ordinary Water And Air.
\end{proof}
\begin{definition}[Satisfies Grazing Water Air Refraction]
  \label{def:physics:phyx-mini-0071:phyxminiproblems-problemphyxmini0071-satisfiesgrazingwaterairrefraction}
  \lean{PhyXMiniProblems.ProblemPhyXMini0071.SatisfiesGrazingWaterAirRefraction}
  \uses{def:physics:phyx-mini-0071:phyxminiproblems-problemphyxmini0071-tankviewingsetup}
  Snell's law for the limiting ray. Because the eye is level with the water surface, the transmitted limiting ray is horizontal and hence makes angle π / 2 with the vertical normal
\end{definition}
\begin{proof}
  This is the declared model definition of Satisfies Grazing Water Air Refraction.
\end{proof}
\begin{definition}[Satisfies Limiting Sightline Geometry]
  \label{def:physics:phyx-mini-0071:phyxminiproblems-problemphyxmini0071-satisfieslimitingsightlinegeometry}
  \lean{PhyXMiniProblems.ProblemPhyXMini0071.SatisfiesLimitingSightlineGeometry}
  \uses{def:physics:phyx-mini-0071:phyxminiproblems-problemphyxmini0071-lengthincentimeters, def:physics:phyx-mini-0071:phyxminiproblems-problemphyxmini0071-tankviewingsetup}
  Right-triangle geometry of the limiting ray across the tank. The angle is measured from the vertical normal, so horizontal run equals vertical drop times its tangent. The depth is additionally constrained to the scale's physical range, without assigning the requested numerical mark
\end{definition}
\begin{proof}
  This is the declared model definition of Satisfies Limiting Sightline Geometry.
\end{proof}
\begin{definition}[Answer Choice]
  \label{def:physics:phyx-mini-0071:phyxminiproblems-problemphyxmini0071-answerchoice}
  \lean{PhyXMiniProblems.ProblemPhyXMini0071.AnswerChoice}
  Labels of the four multiple-choice answers
\end{definition}
\begin{proof}
  This is the declared model definition of Answer Choice.
\end{proof}
\begin{definition}[answer Mark]
  \label{def:physics:phyx-mini-0071:phyxminiproblems-problemphyxmini0071-answermark}
  \lean{PhyXMiniProblems.ProblemPhyXMini0071.answerMark}
  \uses{def:physics:phyx-mini-0071:phyxminiproblems-problemphyxmini0071-depthmark, def:physics:phyx-mini-0071:phyxminiproblems-problemphyxmini0071-answerchoice}
  The depth mark printed beside each answer choice
\end{definition}
\begin{proof}
  This is the declared model definition of answer Mark.
\end{proof}
\begin{definition}[Is Nearest Displayed Mark]
  \label{def:physics:phyx-mini-0071:phyxminiproblems-problemphyxmini0071-isnearestdisplayedmark}
  \lean{PhyXMiniProblems.ProblemPhyXMini0071.IsNearestDisplayedMark}
  \uses{def:physics:phyx-mini-0071:phyxminiproblems-problemphyxmini0071-lengthincentimeters, def:physics:phyx-mini-0071:phyxminiproblems-problemphyxmini0071-depthmark, def:physics:phyx-mini-0071:phyxminiproblems-problemphyxmini0071-depthmarkcentimeters, def:physics:phyx-mini-0071:phyxminiproblems-problemphyxmini0071-tankviewingsetup}
  A limiting sightline selects a displayed mark when its intersection depth is within half of the 10 cm scale spacing of that mark
\end{definition}
\begin{proof}
  This is the declared model definition of Is Nearest Displayed Mark.
\end{proof}
\begin{definition}[Is Correct Answer]
  \label{def:physics:phyx-mini-0071:phyxminiproblems-problemphyxmini0071-iscorrectanswer}
  \lean{PhyXMiniProblems.ProblemPhyXMini0071.IsCorrectAnswer}
  \uses{def:physics:phyx-mini-0071:phyxminiproblems-problemphyxmini0071-tankviewingsetup, def:physics:phyx-mini-0071:phyxminiproblems-problemphyxmini0071-answerchoice, def:physics:phyx-mini-0071:phyxminiproblems-problemphyxmini0071-answermark, def:physics:phyx-mini-0071:phyxminiproblems-problemphyxmini0071-isnearestdisplayedmark}
  A multiple-choice answer names the marked depth selected by the sightline
\end{definition}
\begin{proof}
  This is the declared model definition of Is Correct Answer.
\end{proof}
\begin{definition}[recorded Answer Choice]
  \label{def:physics:phyx-mini-0071:phyxminiproblems-problemphyxmini0071-recordedanswerchoice}
  \lean{PhyXMiniProblems.ProblemPhyXMini0071.recordedAnswerChoice}
  \uses{def:physics:phyx-mini-0071:phyxminiproblems-problemphyxmini0071-answerchoice}
  Dataset metadata: the recorded answer label. It is not a theorem premise
\end{definition}
\begin{proof}
  This is the declared model definition of recorded Answer Choice.
\end{proof}
\begin{lemma}[limiting Sightline Depth formula]
  \label{lem:physics:phyx-mini-0071:phyxminiproblems-problemphyxmini0071-limitingsightlinedepth-formula}
  \lean{PhyXMiniProblems.ProblemPhyXMini0071.limitingSightlineDepth_formula}
  \uses{def:physics:phyx-mini-0071:phyxminiproblems-problemphyxmini0071-lengthincentimeters, def:physics:phyx-mini-0071:phyxminiproblems-problemphyxmini0071-tankviewingsetup, def:physics:phyx-mini-0071:phyxminiproblems-problemphyxmini0071-satisfiesgrazingwaterairrefraction, def:physics:phyx-mini-0071:phyxminiproblems-problemphyxmini0071-satisfieslimitingsightlinegeometry}
  The right-triangle relation and physical acute-angle branch give the limiting depth as tank width divided by the tangent of the underwater incidence angle
\end{lemma}
\begin{proof}
  The conclusion follows from the governing relations and figure readouts named in the statement of limiting Sightline Depth formula.
\end{proof}
\begin{lemma}[limiting Sightline selects sixty centimeter mark]
  \label{lem:physics:phyx-mini-0071:phyxminiproblems-problemphyxmini0071-limitingsightline-selects-sixty-centimeter-mark}
  \lean{PhyXMiniProblems.ProblemPhyXMini0071.limitingSightline_selects_sixty_centimeter_mark}
  \uses{def:physics:phyx-mini-0071:phyxminiproblems-problemphyxmini0071-tankviewingsetup, def:physics:phyx-mini-0071:phyxminiproblems-problemphyxmini0071-matchestankproblemandfigure, def:physics:phyx-mini-0071:phyxminiproblems-problemphyxmini0071-modelsordinarywaterandair, def:physics:phyx-mini-0071:phyxminiproblems-problemphyxmini0071-satisfiesgrazingwaterairrefraction, def:physics:phyx-mini-0071:phyxminiproblems-problemphyxmini0071-satisfieslimitingsightlinegeometry, def:physics:phyx-mini-0071:phyxminiproblems-problemphyxmini0071-isnearestdisplayedmark}
  The 65 cm width, ordinary water/air indices, grazing Snell law, and scale spacing make the limiting depth closest to the 60 cm mark
\end{lemma}
\begin{proof}
  The conclusion follows from the governing relations and figure readouts named in the statement of limiting Sightline selects sixty centimeter mark.
\end{proof}
% --- Archon declaration topology end ---
% --- Archon physics formalization source end ---
